\documentclass[10pt,a4paper]{article}
\usepackage{fontspec,amsmath,amssymb,graphicx}
\usepackage[margin=18mm]{geometry}
\setmainfont{Times New Roman}
\setlength{\parindent}{0pt}
\setlength{\parskip}{5pt}
\allowdisplaybreaks
\sloppy
\emergencystretch=2em
\begin{document}
\section*{BT\_5\_6\_3\_05 — Lời giải bổ sung}
(a) Tích phân trên chu kỳ cho $\lambda\int|\phi|^2=\int|\nabla\phi|^2\ge0$.
(b) Cơ sở đầy đủ là $e^{i(nx+my)}$, $n,m\in\mathbb Z$, $\lambda=n^2+m^2$; phần thực là cos/sin của $nx+my$. Chỉ lấy $n,m\ge0$ như đề sẽ bỏ các mode $nx-my$.
(c) Tích phân trực giao cần tích hai hàm (hoặc liên hợp) và hệ số chuẩn hóa, không phải $\int\phi_{nm}$ đơn lẻ như dòng đề. Cơ sở chuẩn hóa $(2\pi)^{-1}e^{i(nx+my)}$ có $\int_{(-\pi,\pi)^2}\phi_{nm}\overline\phi_{n'm'}=\delta_{nn'}\delta_{mm'}$. Trên $(0,\pi)^2$, các sóng chéo đã in không thỏa công thức trực giao ấy nói chung.
(d) Fourier của $x$ trên $(-\pi,\pi)$ là $2\sum_{n\ge1}(-1)^{n+1}\sin nx/n$. Vì $xy$ là tích hai chuỗi,
\[u=4\sum_{n,m\ge1}\frac{(-1)^{n+m}}{nm}e^{-(n^2+m^2)t}\sin nx\sin my.\]
Giới hạn 0, nghiệm tuần hoàn trơn khi $t>0$; dữ liệu ban đầu được tuần hoàn hóa và có gián đoạn ở mép.\end{document}